\chapter{Evaluation and discussion}
\label{chap:discussion}

In this chapter, we present the results achieved through this work and provide a
comparative analysis of our Jasmin language server implementation against the
existing implementation, focusing on the features each supports.

\section{Evaluation}

While the language server is still under active development, it implements 
core functionality that is present in other language server, and it covers 
basic use-cases and requirements posed in Section~\ref{chap:met:requierements}.

Table~\ref{table:feature-comparison} below illustrates the feature comparison
between the existing Jasmin language server and our implementation. This
comparison highlights the enhancements and capabilities we have introduced,
showcasing the strides made towards improving the Jasmin development
environment.
    
\begin{longtable}[h]{c|c|c}
\caption[Language server features comparison table]{Feature comparison} \label{table:feature-comparison} \\
\hline
Feature Compatibility     & Existing Implementation & Our Work              \\
\hline
\endfirsthead
Diagnostics Support       & Yes                                & Yes                   \\
Goto Definition Support   & No                                 & Yes                   \\
Incremental Parsing       & Yes                                & No                    \\
Symbol Highlight Support  & Yes                                & Yes                   \\
Syntax/Semantic Highlight & No                                 & Syntax Highlight Only \\
Symbol Type on Hover      & No                                 & Yes                  
\end{longtable}

Our language server implementation surpasses the existing solution by
incorporating additional features. Both versions offer essential diagnostics
support; however, we have expanded functionality with the `goto definition'
feature, making code navigation more straightforward for developers.

Additionally, our server displays symbol types on hover, which allows the 
developers to easily get the symbol type.

While our implementation currently supports only syntax highlighting, it sets the
stage for future semantic highlighting enhancements, which will further improve
code analysis and editing capabilities.

\section{Discussion and reflection}

Speaking of the challenges faced during the development of the Jasmin
language server, one significant issue was the niche status of the Jasmin
language itself. Being a relatively specialized and rapidly evolving programming
language, the Jasmin compiler presented several undocumented features. This lack
of comprehensive documentation posed considerable difficulties in understanding
and integrating these features into the language server. 

Another challenge during the implementation was the poor choice of testing 
methodology, in future releases of the language server we expect to add end-to-end 
tests, for this the GitHub actions may suite well.
